%\addcontentsline{toc}{chapter}{Abstract}

\begin{abstract}
\pagenumbering{gobble}

This thesis explores the current state of Automatic Speech Recognition (ASR) with a focus on extending it into areas that are currently not envisioned, namely the temporal dynamics and prosodic dimensions characteristic of spoken language. The established purpose of ASR is the transformation of spoken language into written language that corresponds to some standard orthography. This is a mapping from a highly individual, continuous system of complex sound waves with rapidly varying basic frequencies, overtones, and amplitudes, into a very constrained, categorical, standardized sequence of letters. While this transformation is made possible by the very nature of human language with its underlying digital dimension of discrete phonemes, it is a dramatic reduction of complexity: When we read a text, we don't "read" the voice of the person, their emotional arousal, the emphasis put on words for the purpose of highlighting, the de-emphasis on words that are treated as given, the changing frequency indicating a question or expressing incredulity, and we don't "read" when the speaker speeds up or slows down or is making a pause.

This is largely a good thing: Writing and Reading have different virtues and purposes from Speaking and Listening. However, there are occasions and purposes for which it would be beneficial to capture some of the wealth of speech in writing. For example, movie subtitles could convey certain aspects of the spoken original. From the other end, text-to-speech systems might profit from customized indications about the prosody of their output to come closer to an optimal representation in the acoustic domain.

In this thesis, we will explore how currently available tools can be harnessed to create a new device, "SpeechPrint," that produces a symbolic, i.e., digital, representation of the dynamic and prosodic features of spoken language. For the purpose of the thesis, the target is quite narrowly defined: It should produce outputs that can serve as suggestions for the linguistic transcription and annotation of spoken language. But we will mention use scenarios that go beyond this particular scientific purpose.

With this goal in mind, we will first examine the multidimensional nature of spoken language and current research in linguistics, including the major tools used, Praat and ELAN. Then we will give an overview of currently available speech recognition and analysis tools, namely, wave2vec, HuBERT, Whisper, WhisperX, the Montreal Forced Aligner, ProsodyPro, and the Parselmouth library. We continue with the construction of SpeechPrint and the particular symbolic representation of prosodic features. This is followed by testing SpeechPrint's prosody performance on a variety of spoken data across different languages, which can serve as benchmarks. We will see that the results are promising, but leave room for improvement. Finally, we will discuss a range of use scenarios for the representation of dynamic and prosodic features in forms of written language, and in other, including artistic, representations of speech.

% ADDED HERE — revised abstract for submission (Lonce's feedback: write this first, nail it, one page max 1.5)

Spoken language carries far more information than its orthographic transcription preserves. Pitch movements signal questions, emphasis, and information structure; syllable duration marks prominence and phrase boundaries; amplitude variations highlight the communicatively most important words. These prosodic dimensions are the subject of a substantial body of linguistic research, yet they remain underrepresented in the output of automated speech recognition systems and difficult to annotate consistently without specialist training. This thesis presents SpeechPrint, a computational pipeline that takes a WAV recording as input and produces a structured Praat TextGrid as output, containing five time-aligned tiers: words, phonemes, syllables, fundamental frequency measurements at vowel nuclei, and a symbolic prosody layer.

The symbolic layer encodes prosodic events using a small set of ASCII-compatible symbols: \texttt{/} for a rising pitch movement within a syllable, \texttt{\textbackslash} for a falling movement, \texttt{{\={}}} (U+203E overline) for a high level relative to neighbouring syllables, \texttt{\_} for a low level, \texttt{*} for a prominent accent, and combinations of these for complex events such as \texttt{*{\={}}//} (accented, high, strongly rising) or \texttt{\_\textbackslash\textbackslash} (low, strongly falling). These symbols were chosen to be intuitive, printable in any text environment, and directly mappable to standard prosodic annotation systems such as ToBI. The labels are assigned using recording-adaptive thresholds derived from the normalised median absolute deviation of pitch movements, making the system robust to outlier syllables that would inflate a standard-deviation estimate.

A central methodological contribution is the investigation and correction of octave errors in fundamental frequency extraction. Praat's default autocorrelation pitch tracker, when applied to the English recordings used in this study, produces 238 frames in a single 30-second excerpt where the reported F0 is approximately half the true value --- an error of exactly 12 semitones. This directly inverts the intra-syllable pitch movement used to assign rising and falling labels, causing a genuinely rising syllable to be labelled as falling or vice versa. Five pitch trackers were evaluated and compared quantitatively: Praat's default autocorrelation method; Praat with an increased octave-jump penalty and post-hoc Xu~(1999) spike correction; the probabilistic YIN algorithm (pYIN) from Librosa; the deep convolutional model CREPE (Kim et al., 2018), accessed via the PyTorch port \texttt{torchcrepe}; and PESTO, a self-supervised tracker trained via transposition-equivariant objectives. On the five German GToBI benchmark sentences, all trackers agree within 2\,Hz and no octave errors are detected. On the English and Daakie recordings, pYIN and CREPE agree on a median near 190--214\,Hz while Praat raw reports approximately half that value. Based on this comparison, two pipeline configurations were adopted: pYIN is the primary tracker for read and studio speech with available ASR support (English, German GToBI); CREPE is used for archival recordings of under-resourced languages without available ASR (Daakie, Cabécar), where its data-driven, non-harmonic approach is more robust to irregular phonation.

The pipeline was evaluated against five German sentences hand-annotated according to the GToBI system, using the human word boundaries as input so that alignment error does not confound the comparison. With the final pYIN + five-optimisation pipeline, two of five sentences are fully correct and three are partial (correct word, correct direction, missing accent marker); no sentence places the accent on the wrong syllable. An intermediate CREPE build achieved one fully correct result but placed the pitch direction correctly on all five sentences, illustrating a trade-off between binary accuracy and interpretive quality. The system was also applied to an English minimal-pairs corpus (305.5\,s, 21 prosodic categories), two DoReCo archival recordings --- Daakie (162.7\,s, Oceanic, Vanuatu) and Cabécar (81.8\,s, Chibchan, Costa Rica) --- providing cross-linguistic coverage across four typologically distinct datasets.

The system is designed to support linguists who are not trained in phonetic software. The TextGrid output is directly readable in Praat without additional installation, the symbolic tier can be inspected visually without acoustic expertise, and the pipeline requires only a WAV file and a language code as input. Possible use scenarios range from field work on under-resourced languages, where automatic pre-annotation can reduce transcription time, to subtitle generation for media, to the resynthesis and manipulation of prosody via text-to-speech systems that accept prosodic control parameters.

\bigskip
\noindent\textbf{Keywords:} prosody annotation; fundamental frequency; pitch extraction; forced alignment; TextGrid; SpeechPrint

% ENDED ADDING

\newpage
\end{abstract}
